\subsubsection{Diffusion experiments}

In this set of diffusion experiments, we evaluate the trade-off between concept erasure effectiveness (measured by CLIP score) and image quality preservation (measured by FID - Fréchet Inception Distance). We plot CLIP scores for erased concepts against FID scores aggregated across unrelated concepts to construct Pareto efficiency frontiers for different steering methods.

For each concept erasure task, we measure:
\begin{itemize}
    \item \textbf{CLIP Score}: Presence of the erased concept in generated images (lower is better for successful erasure)
    \item \textbf{FID}: Image quality discrepancy between original and steered images, averaged across unrelated test concepts (lower is better)
\end{itemize}

The Pareto frontier identifies method-beta combinations that cannot be improved in both metrics simultaneously. Results are presented for the available experimental configurations:

\begin{itemize}
    \item Baseline (No Renormalization, No Clipping): \ref{fig:diffusion_pareto_frontier_chart_all_strengths}
    \item With Renormalization, No Clipping: \ref{fig:diffusion_pareto_frontier_chart_all_strengths_renorm}
    \item With Clipping, No Renormalization: \ref{fig:diffusion_pareto_frontier_chart_all_strengths_clip}
\end{itemize}

Note: Each chart includes all available concept erasure tasks for the respective experimental configuration. Tasks with insufficient data (empty files or only baseline methods) are automatically excluded.

